% Options for packages loaded elsewhere
\PassOptionsToPackage{unicode}{hyperref}
\PassOptionsToPackage{hyphens}{url}
%
\documentclass[
  ignorenonframetext,
  aspectratio=169,
]{beamer}
\usepackage{pgfpages}
\setbeamertemplate{caption}[numbered]
\setbeamertemplate{caption label separator}{: }
\setbeamercolor{caption name}{fg=normal text.fg}
\beamertemplatenavigationsymbolshorizontal
% Prevent slide breaks in the middle of a paragraph
\widowpenalties 1 10000
\raggedbottom
\setbeamertemplate{part page}{
  \centering
  \begin{beamercolorbox}[sep=16pt,center]{part title}
    \usebeamerfont{part title}\insertpart\par
  \end{beamercolorbox}
}
\setbeamertemplate{section page}{
  \centering
  \begin{beamercolorbox}[sep=12pt,center]{section title}
    \usebeamerfont{section title}\insertsection\par
  \end{beamercolorbox}
}
\setbeamertemplate{subsection page}{
  \centering
  \begin{beamercolorbox}[sep=8pt,center]{subsection title}
    \usebeamerfont{subsection title}\insertsubsection\par
  \end{beamercolorbox}
}
\AtBeginPart{
  \frame{\partpage}
}
\AtBeginSection{
  \ifbibliography
  \else
    \frame{\sectionpage}
  \fi
}
\AtBeginSubsection{
  \frame{\subsectionpage}
}

\usepackage{amsmath,amssymb}
\usepackage{iftex}
\ifPDFTeX
  \usepackage[T1]{fontenc}
  \usepackage[utf8]{inputenc}
  \usepackage{textcomp} % provide euro and other symbols
\else % if luatex or xetex
  \usepackage{unicode-math}
  \defaultfontfeatures{Scale=MatchLowercase}
  \defaultfontfeatures[\rmfamily]{Ligatures=TeX,Scale=1}
\fi
\usepackage{lmodern}
\usetheme[]{Madrid}
\usecolortheme{dove}
\ifPDFTeX\else  
    % xetex/luatex font selection
\fi
% Use upquote if available, for straight quotes in verbatim environments
\IfFileExists{upquote.sty}{\usepackage{upquote}}{}
\IfFileExists{microtype.sty}{% use microtype if available
  \usepackage[]{microtype}
  \UseMicrotypeSet[protrusion]{basicmath} % disable protrusion for tt fonts
}{}
\usepackage{xcolor}
\newif\ifbibliography
\setlength{\emergencystretch}{3em} % prevent overfull lines
\setcounter{secnumdepth}{-\maxdimen} % remove section numbering


\providecommand{\tightlist}{%
  \setlength{\itemsep}{0pt}\setlength{\parskip}{0pt}}\usepackage{longtable,booktabs,array}
\usepackage{calc} % for calculating minipage widths
\usepackage{caption}
% Make caption package work with longtable
\makeatletter
\def\fnum@table{\tablename~\thetable}
\makeatother
\usepackage{graphicx}
\makeatletter
\newsavebox\pandoc@box
\newcommand*\pandocbounded[1]{% scales image to fit in text height/width
  \sbox\pandoc@box{#1}%
  \Gscale@div\@tempa{\textheight}{\dimexpr\ht\pandoc@box+\dp\pandoc@box\relax}%
  \Gscale@div\@tempb{\linewidth}{\wd\pandoc@box}%
  \ifdim\@tempb\p@<\@tempa\p@\let\@tempa\@tempb\fi% select the smaller of both
  \ifdim\@tempa\p@<\p@\scalebox{\@tempa}{\usebox\pandoc@box}%
  \else\usebox{\pandoc@box}%
  \fi%
}
% Set default figure placement to htbp
\def\fps@figure{htbp}
\makeatother
% definitions for citeproc citations
\NewDocumentCommand\citeproctext{}{}
\NewDocumentCommand\citeproc{mm}{%
  \begingroup\def\citeproctext{#2}\cite{#1}\endgroup}
\makeatletter
 % allow citations to break across lines
 \let\@cite@ofmt\@firstofone
 % avoid brackets around text for \cite:
 \def\@biblabel#1{}
 \def\@cite#1#2{{#1\if@tempswa , #2\fi}}
\makeatother
\newlength{\cslhangindent}
\setlength{\cslhangindent}{1.5em}
\newlength{\csllabelwidth}
\setlength{\csllabelwidth}{3em}
\newenvironment{CSLReferences}[2] % #1 hanging-indent, #2 entry-spacing
 {\begin{list}{}{%
  \setlength{\itemindent}{0pt}
  \setlength{\leftmargin}{0pt}
  \setlength{\parsep}{0pt}
  % turn on hanging indent if param 1 is 1
  \ifodd #1
   \setlength{\leftmargin}{\cslhangindent}
   \setlength{\itemindent}{-1\cslhangindent}
  \fi
  % set entry spacing
  \setlength{\itemsep}{#2\baselineskip}}}
 {\end{list}}
\usepackage{calc}
\newcommand{\CSLBlock}[1]{\hfill\break\parbox[t]{\linewidth}{\strut\ignorespaces#1\strut}}
\newcommand{\CSLLeftMargin}[1]{\parbox[t]{\csllabelwidth}{\strut#1\strut}}
\newcommand{\CSLRightInline}[1]{\parbox[t]{\linewidth - \csllabelwidth}{\strut#1\strut}}
\newcommand{\CSLIndent}[1]{\hspace{\cslhangindent}#1}

\IfFileExists{plex-otf.sty}{
  %% Full TeXlive
  % \usepackage[%
  %   % math,
  %   RM={Scale=0.94},SS={Scale=0.94},SScon={Scale=0.94},TT={Scale=MatchLowercase,FakeStretch=0.9},DefaultFeatures={Ligatures=Common}
  % ]{plex-otf}
}{
  %% TinyTeX
  \usepackage{libertine}
}

%%% Load theme
% https://deic.uab.cat/~iblanes/beamer_gallery/
\IfFileExists{beamerthemegotham.sty}{
  %% Full TeXlive
  \usetheme{gotham}
  \gothamset{
    numbering=totalpagenumber,
    parttocframe default=off,
    sectiontocframe default=off,
    subsectiontocframe default=off,
  }
}{
  %% TinyTeX
  \usetheme{Madrid}
}
\makeatletter
\@ifpackageloaded{caption}{}{\usepackage{caption}}
\AtBeginDocument{%
\ifdefined\contentsname
  \renewcommand*\contentsname{Table of contents}
\else
  \newcommand\contentsname{Table of contents}
\fi
\ifdefined\listfigurename
  \renewcommand*\listfigurename{Список иллюстраций}
\else
  \newcommand\listfigurename{Список иллюстраций}
\fi
\ifdefined\listtablename
  \renewcommand*\listtablename{Список таблиц}
\else
  \newcommand\listtablename{Список таблиц}
\fi
\ifdefined\figurename
  \renewcommand*\figurename{Figure}
\else
  \newcommand\figurename{Figure}
\fi
\ifdefined\tablename
  \renewcommand*\tablename{Table}
\else
  \newcommand\tablename{Table}
\fi
}
\@ifpackageloaded{float}{}{\usepackage{float}}
\floatstyle{ruled}
\@ifundefined{c@chapter}{\newfloat{codelisting}{h}{lop}}{\newfloat{codelisting}{h}{lop}[chapter]}
\floatname{codelisting}{Listing}
\newcommand*\listoflistings{\listof{codelisting}{Листинги}}
\makeatother
\makeatletter
\makeatother
\makeatletter
\@ifpackageloaded{caption}{}{\usepackage{caption}}
\@ifpackageloaded{subcaption}{}{\usepackage{subcaption}}
\makeatother

\usepackage{csquotes}
\usepackage{bookmark}

\IfFileExists{xurl.sty}{\usepackage{xurl}}{} % add URL line breaks if available
\urlstyle{same} % disable monospaced font for URLs
\hypersetup{
  pdftitle={Лабораторная работа № 6},
  pdfauthor={Гайдук Софья Сергеевна},
  hidelinks,
  pdfcreator={LaTeX via pandoc}}


\title{Лабораторная работа № 6}
\subtitle{Презентация}
\author{Гайдук Софья Сергеевна}
\date{2026-03-20}

\begin{document}
\frame{\titlepage}

\renewcommand*\contentsname{Содержание}
\begin{frame}[allowframebreaks]
  \frametitle{Содержание}
  \setcounter{tocdepth}{1}
  \tableofcontents
\end{frame}

\begin{frame}{Title}
\phantomsection\label{title}
\end{frame}

\section{Информация}\label{ux438ux43dux444ux43eux440ux43cux430ux446ux438ux44f}

\begin{frame}{Докладчик}
\phantomsection\label{ux434ux43eux43aux43bux430ux434ux447ux438ux43a}
\begin{columns}[c]
\begin{column}{0.7\linewidth}
\begin{itemize}[<+->]
\tightlist
\item
  Гайдук Софья Сергеевна
\item
  студент
\item
  Российский университет дружбы народов им. П. Лумумбы
\item
  \href{mailto:1032253645@rudn.ru}{\nolinkurl{1032253645@rudn.ru}}
\item
  \url{https://github.com/SofiaGayduk/study_2025-2026_os-intro}
\end{itemize}
\end{column}

\begin{column}{0.3\linewidth}
\end{column}
\end{columns}
\end{frame}

\begin{frame}{Цель работы}
\phantomsection\label{ux446ux435ux43bux44c-ux440ux430ux431ux43eux442ux44b}
Приобретение практических навыков взаимодействия пользователя с системой
посредством командной строки.
\end{frame}

\begin{frame}{Выполнение лабораторной работы}
\phantomsection\label{ux432ux44bux43fux43eux43bux43dux435ux43dux438ux435-ux43bux430ux431ux43eux440ux430ux442ux43eux440ux43dux43eux439-ux440ux430ux431ux43eux442ux44b}
\end{frame}

\begin{frame}{Определение имени}
\phantomsection\label{ux43eux43fux440ux435ux434ux435ux43bux435ux43dux438ux435-ux438ux43cux435ux43dux438}
Определим полное имя нашего домашнего каталога ({[}рис. 1{]}).

\begin{figure}

\centering{

\includegraphics[width=0.7\linewidth,height=\textheight,keepaspectratio]{image/image1.jpg}

}

\caption{\label{fig-001}image1}

\end{figure}%
\end{frame}

\begin{frame}{Команда ls}
\phantomsection\label{ux43aux43eux43cux430ux43dux434ux430-ls}
Перейдем в каталог /tmp ({[}рис. 2{]}).

\begin{figure}

\centering{

\includegraphics[width=0.7\linewidth,height=\textheight,keepaspectratio]{image/image2.jpg}

}

\caption{\label{fig-002}image2}

\end{figure}%
\end{frame}

\begin{frame}{Команда ls}
\phantomsection\label{ux43aux43eux43cux430ux43dux434ux430-ls-1}
Выведем на экран содержимое каталога /tmp командой ls ({[}рис. 3{]}).

\begin{figure}

\centering{

\includegraphics[width=0.7\linewidth,height=\textheight,keepaspectratio]{image/image3.jpg}

}

\caption{\label{fig-003}image3}

\end{figure}%
\end{frame}

\begin{frame}{Команда ls}
\phantomsection\label{ux43aux43eux43cux430ux43dux434ux430-ls-2}
Выведем на экран содержимое каталога /tmp командой ls с различными
опциями. ({[}рис. 4{]}).

\begin{figure}

\centering{

\includegraphics[width=0.7\linewidth,height=\textheight,keepaspectratio]{image/image4.jpg}

}

\caption{\label{fig-004}image4}

\end{figure}%
\end{frame}

\begin{frame}{Команда ls}
\phantomsection\label{ux43aux43eux43cux430ux43dux434ux430-ls-3}
Выведем на экран содержимое каталога /tmp командой ls с различными
опциями. ({[}рис. 5{]}).

\begin{figure}

\centering{

\includegraphics[width=0.7\linewidth,height=\textheight,keepaspectratio]{image/image5.jpg}

}

\caption{\label{fig-005}image5}

\end{figure}%
\end{frame}

\begin{frame}{Просмотр катологов}
\phantomsection\label{ux43fux440ux43eux441ux43cux43eux442ux440-ux43aux430ux442ux43eux43bux43eux433ux43eux432}
Определим, есть ли в каталоге /var/spool подкаталог с именем cron
({[}рис. 6{]}).

\begin{figure}

\centering{

\includegraphics[width=0.7\linewidth,height=\textheight,keepaspectratio]{image/image6.jpg}

}

\caption{\label{fig-006}image6}

\end{figure}%
\end{frame}

\begin{frame}{Просмотр катологов}
\phantomsection\label{ux43fux440ux43eux441ux43cux43eux442ux440-ux43aux430ux442ux43eux43bux43eux433ux43eux432-1}
Перейдем в домашний каталог и выведем на экран содержимое. Определим
владельца файлов и подкаталогов ({[}рис. 7{]}).

\begin{figure}

\centering{

\includegraphics[width=0.7\linewidth,height=\textheight,keepaspectratio]{image/image7.jpg}

}

\caption{\label{fig-007}image7}

\end{figure}%
\end{frame}

\begin{frame}{Создание каталогов}
\phantomsection\label{ux441ux43eux437ux434ux430ux43dux438ux435-ux43aux430ux442ux430ux43bux43eux433ux43eux432}
Создадим новый каталог newdir ({[}рис. 8{]}).

\begin{figure}

\centering{

\includegraphics[width=0.7\linewidth,height=\textheight,keepaspectratio]{image/image8.jpg}

}

\caption{\label{fig-008}image8}

\end{figure}%

В каталоге \textasciitilde/newdir создадим новый каталог с именем
morefun ({[}рис. 9{]}).

\begin{figure}

\centering{

\includegraphics[width=0.7\linewidth,height=\textheight,keepaspectratio]{image/image9.jpg}

}

\caption{\label{fig-009}image9}

\end{figure}%
\end{frame}

\begin{frame}{Создание и удаление каталогов}
\phantomsection\label{ux441ux43eux437ux434ux430ux43dux438ux435-ux438-ux443ux434ux430ux43bux435ux43dux438ux435-ux43aux430ux442ux430ux43bux43eux433ux43eux432}
Создадим одной командой три новых каталога: letters, memos, misk. Затем
удалим эти каталоги одной командой ({[}рис. 10{]}, {[}рис. 11{]}).

\begin{figure}

\centering{

\includegraphics[width=0.7\linewidth,height=\textheight,keepaspectratio]{image/image10.jpg}

}

\caption{\label{fig-010}image10}

\end{figure}%

\begin{figure}

\centering{

\includegraphics[width=0.7\linewidth,height=\textheight,keepaspectratio]{image/image11.jpg}

}

\caption{\label{fig-011}image11}

\end{figure}%
\end{frame}

\begin{frame}{Удаление каталогов}
\phantomsection\label{ux443ux434ux430ux43bux435ux43dux438ux435-ux43aux430ux442ux430ux43bux43eux433ux43eux432}
Попробуем удалить ранее созданный каталог \textasciitilde/newdir
командой rm ({[}рис. 12{]}).

\begin{figure}

\centering{

\includegraphics[width=0.7\linewidth,height=\textheight,keepaspectratio]{image/image12.jpg}

}

\caption{\label{fig-012}image12}

\end{figure}%

Удалим каталог \textasciitilde/newdir/morefun из домашнего каталога
({[}рис. 13{]}, {[}рис. 14{]}).

\begin{figure}

\centering{

\includegraphics[width=0.7\linewidth,height=\textheight,keepaspectratio]{image/image13.jpg}

}

\caption{\label{fig-013}image13}

\end{figure}%

\begin{figure}

\centering{

\includegraphics[width=0.7\linewidth,height=\textheight,keepaspectratio]{image/image14.jpg}

}

\caption{\label{fig-014}image14}

\end{figure}%
\end{frame}

\begin{frame}{Команда man}
\phantomsection\label{ux43aux43eux43cux430ux43dux434ux430-man}
С помощью команды man определим, какую опцию команды ls нужно
использовать для просмотра содержимое не только указанного каталога, но
и подкаталогов, входящих в него ({[}рис. 15{]}, {[}рис. 16{]}).

\begin{figure}

\centering{

\includegraphics[width=0.7\linewidth,height=\textheight,keepaspectratio]{image/image15.jpg}

}

\caption{\label{fig-015}image15}

\end{figure}%

\begin{figure}

\centering{

\includegraphics[width=0.7\linewidth,height=\textheight,keepaspectratio]{image/image16.jpg}

}

\caption{\label{fig-016}image16}

\end{figure}%
\end{frame}

\begin{frame}{Команда man}
\phantomsection\label{ux43aux43eux43cux430ux43dux434ux430-man-1}
С помощью команды man определим набор опций команды ls, позволяющий
отсортировать по времени последнего изменения выводимый список
содержимого каталога с развёрнутым описанием файлов ({[}рис. 17{]},
{[}рис. 18{]}).

\begin{figure}

\centering{

\includegraphics[width=0.7\linewidth,height=\textheight,keepaspectratio]{image/image17.jpg}

}

\caption{\label{fig-017}image17}

\end{figure}%

\begin{figure}

\centering{

\includegraphics[width=0.7\linewidth,height=\textheight,keepaspectratio]{image/image18.jpg}

}

\caption{\label{fig-018}image18}

\end{figure}%
\end{frame}

\begin{frame}{Команда man}
\phantomsection\label{ux43aux43eux43cux430ux43dux434ux430-man-2}
Проверим эти команды ({[}рис. 19{]}, {[}рис. 20{]}).

\begin{figure}

\centering{

\includegraphics[width=0.7\linewidth,height=\textheight,keepaspectratio]{image/image19.jpg}

}

\caption{\label{fig-019}image19}

\end{figure}%

\begin{figure}

\centering{

\includegraphics[width=0.7\linewidth,height=\textheight,keepaspectratio]{image/image20.jpg}

}

\caption{\label{fig-020}image20}

\end{figure}%
\end{frame}

\begin{frame}{Команда man}
\phantomsection\label{ux43aux43eux43cux430ux43dux434ux430-man-3}
Используем команду man для просмотра описания следующих команд: cd, pwd,
mkdir, rmdir, rm. Основные опции этих команд ({[}рис. 21{]}).

\begin{figure}

\centering{

\includegraphics[width=0.7\linewidth,height=\textheight,keepaspectratio]{image/image21.jpg}

}

\caption{\label{fig-021}image21}

\end{figure}%
\end{frame}

\begin{frame}{Команда history}
\phantomsection\label{ux43aux43eux43cux430ux43dux434ux430-history}
Используя информацию, полученную при помощи команды history, выполним
модификацию и исполнение нескольких команд из буфера команд ({[}рис.
22{]}, {[}рис. 23{]}).

\begin{figure}

\centering{

\includegraphics[width=0.7\linewidth,height=\textheight,keepaspectratio]{image/image22.jpg}

}

\caption{\label{fig-022}image22}

\end{figure}%

\begin{figure}

\centering{

\includegraphics[width=0.7\linewidth,height=\textheight,keepaspectratio]{image/image23.jpg}

}

\caption{\label{fig-023}image23}

\end{figure}%
\end{frame}

\begin{frame}{Выводы}
\phantomsection\label{ux432ux44bux432ux43eux434ux44b}
Мы приобрели практические навыки взаимодействия пользователя с системой
посредством командной строки.
\end{frame}

\begin{frame}{Список литературы}
\phantomsection\label{ux441ux43fux438ux441ux43eux43a-ux43bux438ux442ux435ux440ux430ux442ux443ux440ux44b}
1.Kulyabov. Лабораторная работа № 6. Основы интерфейса взаимодействия
пользователя с системой Unix на уровне командной строки. RUDN
\end{frame}

\begin{frame}{Список литературы}
\phantomsection\label{ux441ux43fux438ux441ux43eux43a-ux43bux438ux442ux435ux440ux430ux442ux443ux440ux44b-1}
1.Kulyabov. Лабораторная работа № 6. Основы интерфейса взаимодействия
пользователя с системой Unix на уровне командной строки. RUDN

\phantomsection\label{refs}
\begin{CSLReferences}{0}{1}
\end{CSLReferences}
\end{frame}




\end{document}
